\chapter{Banque d'exercices et de problèmes}

\newcommand{\calcoui}{\textcolor{VE!70!black}{\textbf{Calculatrice : autorisée}}}
\newcommand{\calcnon}{\textcolor{OR!80!black}{\textbf{Calculatrice : non nécessaire}}}
\renewcommand{\corrige}{\TeacherOnly{\enspace\textcolor{OR!75}{\small$\triangleright$\,\textit{corrigé}}}}
\renewcommand{\corrigeP}{\TeacherOnly{\enspace\textcolor{OR!75}{\small$\triangleright$\,\textit{corrigé}}}}

\begin{tcolorbox}[enhanced,breakable,frame hidden,
  borderline west={2pt}{0pt}{BN},colback=BN!3,sharp corners,
  left=12pt,right=10pt,top=7pt,bottom=7pt]
Chaque bloc est autonome et peut être inséré directement dans une évaluation.
L'indication relative à la calculatrice s'applique à l'ensemble du bloc.
Les exercices sont classés par domaine et par niveau de difficulté.
\end{tcolorbox}

\section{Repère diagnostique}

\begin{exo}[Calculs et méthodes de base \calcnon \corrige]{1}
\begin{enumerate}
  \item Calculer et donner le résultat sous forme irréductible :
  \[
    A=\frac34-\frac56
    \qquad\text{et}\qquad
    B=\frac{5}{12}+\frac78.
  \]
  \item Simplifier $2^3\times2^{-5}$ et $\sqrt{81}$.
  \item Développer $(x+3)(x-2)$, puis factoriser $4x^2-12x$.
  \item Résoudre l'équation $7x-5=16$.
\end{enumerate}
\end{exo}

\begin{exo}[Repérage, vecteurs et données \calcnon \corrige]{1}
Dans un repère orthonormé, on donne $A(-1;2)$ et $B(3;5)$.
\begin{enumerate}
  \item Déterminer les coordonnées du vecteur $\vv{AB}$.
  \item Calculer les coordonnées du milieu $I$ de $[AB]$.
  \item Calculer la longueur $AB$.
  \item Pour la série $6,8,8,10,13$, déterminer la moyenne, la médiane
  et l'étendue.
\end{enumerate}
\end{exo}

\section{Logique et algorithmique}

\begin{exo}[Propositions et valeurs de vérité \calcnon \corrige]{1}
On considère les phrases suivantes :
\[
\begin{aligned}
  P&:\text{ « $11$ est un nombre premier »},\\
  Q&:\text{ « $x+4=9$ »},\\
  R&:\text{ « pour tout $x\in\R$, $x^2+1>0$ »}.
\end{aligned}
\]
\begin{enumerate}
  \item Indiquer celles qui sont des propositions.
  \item Donner la valeur de vérité de chaque proposition.
  \item Transformer la phrase $Q$ en proposition vraie en précisant une
  valeur de $x$.
\end{enumerate}
\end{exo}

\begin{exo}[Connecteurs logiques \calcnon \corrige]{1}
On pose $P$ : « $n$ est pair » et $Q$ : « $n$ est multiple de $3$ ».
\begin{enumerate}
  \item Traduire en français $P\land Q$ puis $P\lor Q$.
  \item Donner un entier vérifiant $P\land Q$.
  \item Donner un entier vérifiant $P\lor Q$ mais pas $P\land Q$.
  \item Écrire la négation de $P\lor Q$.
\end{enumerate}
\end{exo}

\begin{exo}[Implication et équivalence \calcnon \corrige]{2}
Pour un entier $n$, on considère :
\[
  P:\text{ « $n$ est multiple de $4$ »},\qquad
  Q:\text{ « $n$ est pair »}.
\]
\begin{enumerate}
  \item Dire si l'implication $P\Rightarrow Q$ est vraie. Justifier.
  \item Dire si l'implication $Q\Rightarrow P$ est vraie. Donner, si
  nécessaire, un contre-exemple.
  \item Les propositions $P$ et $Q$ sont-elles équivalentes ?
  \item Écrire la contraposée de $P\Rightarrow Q$.
\end{enumerate}
\end{exo}

\begin{exo}[Quantificateurs et négations \calcnon \corrige]{2}
\begin{enumerate}
  \item Écrire avec des quantificateurs : « tout réel possède un carré
  positif ou nul ».
  \item Écrire avec des quantificateurs : « l'équation $x^2=7$ possède
  une solution réelle ».
  \item Nier chacune des deux propositions précédentes.
  \item Déterminer la valeur de vérité de chaque proposition.
\end{enumerate}
\end{exo}

\begin{exo}[Exécution d'un algorithme conditionnel \calcnon \corrige]{1}
On applique l'algorithme suivant à un réel $x$ :
\[
  y\leftarrow 2x+1;\qquad
  \text{si }y\geq0\text{ alors }z\leftarrow y
  \text{ sinon }z\leftarrow-y.
\]
\begin{enumerate}
  \item Compléter un tableau donnant les valeurs de $y$ et $z$ pour
  $x=-3$, $x=-\frac12$ et $x=2$.
  \item Quelle expression mathématique l'algorithme calcule-t-il ?
  \item Modifier la condition afin que l'algorithme affiche $0$ lorsque
  $y=0$ et distingue strictement les deux autres cas.
\end{enumerate}
\end{exo}

\begin{exo}[Boucle et cumul \calcnon \corrige]{2}
On initialise $S\leftarrow0$, puis on répète pour $k$ allant de $1$ à $5$ :
\[
  S\leftarrow S+2k.
\]
\begin{enumerate}
  \item Recopier et compléter un tableau donnant $k$ et $S$ après chaque
  passage dans la boucle.
  \item Donner la valeur finale de $S$.
  \item Écrire une version qui calcule la somme des $n$ premiers nombres
  pairs positifs.
  \item Adapter l'algorithme pour additionner les nombres impairs de $1$ à $9$.
\end{enumerate}
\end{exo}

\begin{exo}[Concevoir un algorithme \calcnon \corrige]{2}
On saisit un entier naturel $n$.
\begin{enumerate}
  \item Écrire un algorithme qui affiche « pair » ou « impair ».
  \item Compléter cet algorithme pour afficher les cinq premiers multiples
  positifs de $n$.
  \item Exécuter l'algorithme obtenu pour $n=6$.
  \item Expliquer le rôle de la variable de boucle.
\end{enumerate}
\end{exo}

\section{Nombres réels et calculs}

\begin{exo}[Ensembles de nombres \calcnon \corrige]{1}
Pour chacun des nombres
\[
 -5,\quad \frac73,\quad 0{,}125,\quad \sqrt{49},\quad \sqrt{2},
\]
\begin{enumerate}
  \item donner le plus petit ensemble parmi $\N$, $\Z$, $\D$, $\Q$ et $\R$
  auquel il appartient ;
  \item ranger les cinq nombres dans l'ordre croissant ;
  \item donner une valeur décimale approchée de $\sqrt2$ au centième.
\end{enumerate}
\end{exo}

\begin{exo}[Fractions et priorités opératoires \calcnon \corrige]{1}
Calculer en détaillant les étapes :
\[
 A=\frac56-\frac14,\qquad
 B=\frac23+\frac57,\qquad
 C=\frac{3-\frac12}{1+\frac14}.
\]
\begin{enumerate}
  \item Donner chaque résultat sous forme irréductible.
  \item Comparer $A$ et $B$ sans utiliser leurs écritures décimales.
\end{enumerate}
\end{exo}

\begin{exo}[Puissances et écriture scientifique \calcnon \corrige]{1}
\begin{enumerate}
  \item Simplifier $2^3\times2^{-5}$ et $\dfrac{3^7}{3^4}$.
  \item Écrire $0{,}000\,072$ en notation scientifique.
  \item Écrire $5{,}4\times10^6$ sous forme décimale.
  \item Calculer
  \[
    D=\frac{(4\times10^3)(3\times10^{-5})}{6\times10^{-2}}
  \]
  et donner le résultat en notation scientifique.
\end{enumerate}
\end{exo}

\begin{exo}[Radicaux \calcnon \corrige]{1}
\begin{enumerate}
  \item Simplifier $\sqrt{48}$, $\sqrt{75}$ et $\sqrt{147}$.
  \item Réduire $3\sqrt{12}-2\sqrt{27}+\sqrt{48}$.
  \item Développer $(\sqrt5+2)(\sqrt5-2)$.
  \item Comparer $3\sqrt2$ et $\sqrt{20}$ en justifiant.
\end{enumerate}
\end{exo}

\begin{exo}[Intervalles et valeur absolue \calcnon \corrige]{1}
\begin{enumerate}
  \item Traduire $-2\leq x<5$ sous forme d'intervalle.
  \item Traduire $x\in]-\infty;3]\cup[7;+\infty[$ sous forme d'inégalités.
  \item Résoudre $|x-1|\leq3$.
  \item Résoudre $|2x+1|>5$.
  \item Représenter les deux ensembles de solutions sur une droite graduée.
\end{enumerate}
\end{exo}

\begin{exo}[Encadrements et valeurs approchées \calcnon \corrige]{2}
\begin{enumerate}
  \item Encadrer $\sqrt{30}$ entre deux entiers consécutifs.
  \item En déduire un encadrement de $2\sqrt{30}-3$ par deux entiers.
  \item Sachant que $1{,}414<\sqrt2<1{,}415$, encadrer $5\sqrt2+1$.
  \item Donner une valeur approchée de $5\sqrt2+1$ au centième.
\end{enumerate}
\end{exo}

\section{Équations, inéquations et polynômes}

\begin{exo}[Trinôme : discriminant et racines \calcnon \corrige]{1}
On considère $f(x)=2x^2-5x-3$.
\begin{enumerate}
  \item Identifier $a$, $b$ et $c$, puis calculer le discriminant.
  \item Résoudre l'équation $f(x)=0$.
  \item Factoriser $f(x)$.
  \item Vérifier les résultats obtenus en développant.
\end{enumerate}
\end{exo}

\begin{exo}[Forme canonique et variations \calcnon \corrige]{2}
On considère $g(x)=x^2-6x+5$.
\begin{enumerate}
  \item Mettre $g$ sous forme canonique.
  \item Donner les coordonnées du sommet de la parabole.
  \item Dresser le tableau de variations de $g$.
  \item Résoudre $g(x)\leq0$.
  \item Retrouver le résultat précédent à l'aide de la forme factorisée.
\end{enumerate}
\end{exo}

\begin{exo}[Étude du signe d'un trinôme \calcnon \corrige]{2}
Soit $h(x)=-2x^2+8x-6$.
\begin{enumerate}
  \item Calculer les racines de $h$.
  \item Factoriser $h(x)$.
  \item Dresser le tableau de signes de $h$.
  \item Résoudre $h(x)>0$, puis $h(x)\leq0$.
\end{enumerate}
\end{exo}

\begin{exo}[Polynôme du troisième degré \calcnon \corrige]{2}
On considère $P(x)=x^3-2x^2-5x+6$.
\begin{enumerate}
  \item Vérifier que $1$ est une racine de $P$.
  \item Déterminer un trinôme $Q$ tel que $P(x)=(x-1)Q(x)$.
  \item Factoriser complètement $P$.
  \item Résoudre $P(x)=0$.
  \item Dresser le tableau de signes de $P$.
\end{enumerate}
\end{exo}

\begin{exo}[Équation rationnelle \calcnon \corrige]{2}
On considère l'équation
\[
  \frac{2x-1}{x-3}=\frac{x+2}{x+1}.
\]
\begin{enumerate}
  \item Déterminer les valeurs interdites.
  \item Transformer l'équation en une équation polynomiale.
  \item Résoudre l'équation obtenue.
  \item Vérifier que les solutions respectent les conditions d'existence.
\end{enumerate}
\end{exo}

\begin{exo}[Inéquation rationnelle \calcnon \corrige]{3}
Résoudre dans $\R$ :
\[
  \frac{x-2}{x+1}\geq0.
\]
\begin{enumerate}
  \item Donner la valeur interdite et la valeur qui annule le numérateur.
  \item Construire le tableau de signes du quotient.
  \item Donner l'ensemble des solutions sous forme d'intervalle.
  \item Contrôler le résultat en testant une valeur dans chaque intervalle.
\end{enumerate}
\end{exo}

\section{Géométrie analytique et vecteurs}

\begin{exo}[Équation cartésienne d'une droite \calcnon \corrige]{1}
Dans un repère, on donne $A(-1;1)$ et $B(3;3)$.
\begin{enumerate}
  \item Calculer les coordonnées de $\vv{AB}$.
  \item Donner un vecteur normal à la droite $(AB)$.
  \item Déterminer une équation cartésienne de $(AB)$.
  \item Vérifier que les points $A$ et $B$ satisfont cette équation.
\end{enumerate}
\end{exo}

\begin{exo}[Formes d'une équation de droite \calcnon \corrige]{2}
On considère la droite $d:2x-y+3=0$.
\begin{enumerate}
  \item Donner un vecteur directeur de $d$.
  \item Écrire une équation réduite de $d$.
  \item Donner une représentation paramétrique de $d$.
  \item Déterminer les coordonnées des intersections de $d$ avec les axes.
\end{enumerate}
\end{exo}

\begin{exo}[Parallélisme et perpendicularité \calcnon \corrige]{2}
On donne $A(2;-1)$ et la droite $d:3x-2y+4=0$.
\begin{enumerate}
  \item Déterminer une équation de la droite $d_1$ passant par $A$ et
  parallèle à $d$.
  \item Déterminer une équation de la droite $d_2$ passant par $A$ et
  perpendiculaire à $d$.
  \item Calculer le point d'intersection de $d_2$ et $d$.
\end{enumerate}
\end{exo}

\begin{exo}[Intersection de deux droites \calcnon \corrige]{1}
On considère
\[
 d_1:2x+y-7=0
 \qquad\text{et}\qquad
 d_2:x-3y+5=0.
\]
\begin{enumerate}
  \item Justifier que $d_1$ et $d_2$ ne sont pas parallèles.
  \item Résoudre le système formé par leurs équations.
  \item En déduire les coordonnées de leur point d'intersection.
  \item Vérifier ce résultat dans les deux équations.
\end{enumerate}
\end{exo}

\begin{exo}[Équation d'un cercle \calcnon \corrige]{1}
On considère
\[
 \mathcal C:(x-2)^2+(y+1)^2=16.
\]
\begin{enumerate}
  \item Donner le centre et le rayon de $\mathcal C$.
  \item Dire si $A(2;3)$ appartient à $\mathcal C$.
  \item Dire si $B(6;2)$ est intérieur ou extérieur au cercle.
  \item Développer l'équation de $\mathcal C$.
\end{enumerate}
\end{exo}

\begin{exo}[Retrouver un cercle \calcnon \corrige]{2}
Le cercle $\mathcal C$ a pour diamètre $[AB]$, avec $A(-2;1)$ et $B(4;5)$.
\begin{enumerate}
  \item Déterminer les coordonnées de son centre $\Omega$.
  \item Calculer son rayon.
  \item Écrire une équation cartésienne de $\mathcal C$.
  \item Étudier l'appartenance du point $M(1;5)$ à $\mathcal C$.
\end{enumerate}
\end{exo}

\begin{exo}[Intersection d'une droite et d'un cercle \calcnon \corrige]{3}
On considère le cercle $\mathcal C:x^2+y^2=25$ et la droite $d:y=3$.
\begin{enumerate}
  \item Résoudre le système formé par les équations de $\mathcal C$ et $d$.
  \item Donner les coordonnées des points d'intersection.
  \item Calculer la longueur de la corde découpée par $d$.
  \item Comparer la distance du centre de $\mathcal C$ à $d$ avec le rayon.
\end{enumerate}
\end{exo}

\begin{exo}[Colinéarité et alignement \calcnon \corrige]{2}
Dans un repère, on donne $A(-2;1)$, $B(1;3)$ et $C(4;5)$.
\begin{enumerate}
  \item Calculer $\vv{AB}$ et $\vv{AC}$.
  \item Montrer que les vecteurs sont colinéaires.
  \item En déduire que $A$, $B$ et $C$ sont alignés.
  \item Déterminer le réel $k$ tel que $\vv{AC}=k\vv{AB}$.
\end{enumerate}
\end{exo}

\begin{exo}[Produit scalaire et nature d'un triangle \calcnon \corrige]{2}
Dans un repère orthonormé, on donne $A(0;1)$, $B(4;3)$ et $C(2;-3)$.
\begin{enumerate}
  \item Calculer $\vv{AB}$ et $\vv{AC}$.
  \item Calculer $\vv{AB}\cdot\vv{AC}$.
  \item En déduire la nature du triangle $ABC$ en $A$.
  \item Calculer $AB$ et $AC$, puis préciser complètement la nature de $ABC$.
\end{enumerate}
\end{exo}

\section{Fonctions et trigonométrie}

\begin{exo}[Domaine et images \calcnon \corrige]{1}
On considère $f(x)=\sqrt{x+2}$ et $g(x)=\dfrac{1}{x-3}$.
\begin{enumerate}
  \item Déterminer $D_f$ et $D_g$.
  \item Calculer $f(2)$ et $g(5)$.
  \item Déterminer les antécédents de $2$ par $f$.
  \item Expliquer pourquoi $3$ n'a pas d'image par $g$.
\end{enumerate}
\end{exo}

\begin{exo}[Fonctions de référence \calcnon \corrige]{1}
\begin{enumerate}
  \item Donner le domaine de définition de $x\mapsto x^2$, $x\mapsto x^3$,
  $x\mapsto|x|$ et $x\mapsto\sqrt{x}$.
  \item Préciser lesquelles sont paires ou impaires.
  \item Donner les variations de chacune sur son domaine.
  \item Résoudre graphiquement ou algébriquement $x^2=4$ et $|x|=3$.
\end{enumerate}
\end{exo}

\begin{exo}[Lecture d'un tableau de variations \calcnon \corrige]{2}
Une fonction $f$ est définie sur $[-4;5]$. Elle est décroissante de $3$ à
$-2$ sur $[-4;1]$, puis croissante de $-2$ à $6$ sur $[1;5]$.
\begin{enumerate}
  \item Dresser le tableau de variations de $f$.
  \item Donner le minimum de $f$ et la valeur où il est atteint.
  \item Comparer $f(-3)$ et $f(0)$.
  \item Peut-on comparer avec certitude $f(0)$ et $f(3)$ ? Justifier.
\end{enumerate}
\end{exo}

\begin{exo}[Fonction valeur absolue \calcnon \corrige]{2}
On considère $f(x)=|x-2|-1$.
\begin{enumerate}
  \item Calculer $f(-1)$, $f(0)$, $f(2)$ et $f(5)$.
  \item Résoudre $f(x)=0$.
  \item Résoudre $f(x)\leq0$.
  \item Donner les variations de $f$.
  \item Tracer sa courbe dans un repère.
\end{enumerate}
\end{exo}

\begin{exo}[Fonction quadratique \calcnon \corrige]{2}
On considère $g(x)=x^2-4x+3$.
\begin{enumerate}
  \item Donner les formes canonique et factorisée de $g$.
  \item Déterminer son sommet et son axe de symétrie.
  \item Dresser son tableau de variations.
  \item Résoudre $g(x)\leq0$.
  \item Tracer l'allure de la courbe en faisant apparaître les points remarquables.
\end{enumerate}
\end{exo}

\begin{exo}[Angles et cercle trigonométrique \calcnon \corrige]{1}
\begin{enumerate}
  \item Convertir $150^\circ$ et $-45^\circ$ en radians.
  \item Convertir $\dfrac{5\pi}{6}$ et $\dfrac{7\pi}{4}$ en degrés.
  \item Placer ces quatre angles sur un cercle trigonométrique.
  \item Donner une mesure principale de $\dfrac{17\pi}{6}$.
\end{enumerate}
\end{exo}

\begin{exo}[Valeurs trigonométriques \calcnon \corrige]{2}
\begin{enumerate}
  \item Donner $\cos(2\pi/3)$ et $\sin(2\pi/3)$.
  \item Donner $\cos(5\pi/6)$ et $\sin(5\pi/6)$.
  \item Sachant que $\sin\theta=\dfrac35$ et que $\theta$ appartient au
  deuxième quadrant, déterminer $\cos\theta$.
  \item En déduire $\tan\theta$.
\end{enumerate}
\end{exo}

\section{Statistique descriptive}

\begin{exo}[Série brute \calcnon \corrige]{1}
On considère la série
\[
  4,\ 6,\ 6,\ 7,\ 8,\ 8,\ 8,\ 10,\ 11.
\]
\begin{enumerate}
  \item Donner l'effectif total et le mode.
  \item Calculer la moyenne.
  \item Déterminer la médiane et l'étendue.
  \item Expliquer l'effet sur la moyenne du remplacement de $11$ par $20$.
\end{enumerate}
\end{exo}

\begin{exo}[Tableau d'effectifs \calcnon \corrige]{2}
Les notes obtenues sont résumées ci-dessous :
\[
\begin{array}{c|ccccc}
\text{Note} & 6 & 8 & 10 & 12 & 14\\ \hline
\text{Effectif} & 2 & 5 & 7 & 4 & 2
\end{array}
\]
\begin{enumerate}
  \item Calculer l'effectif total.
  \item Calculer les fréquences, en pourcentage.
  \item Déterminer la moyenne et la médiane.
  \item Construire le diagramme en bâtons correspondant.
\end{enumerate}
\end{exo}

\begin{exo}[Variance et écart-type \calcoui \corrige]{2}
On considère la série $0,1,2,2,3,4,4,8$.
\begin{enumerate}
  \item Calculer la moyenne et la médiane.
  \item Calculer la variance.
  \item En déduire l'écart-type, arrondi au centième.
  \item Reprendre les calculs après suppression de la valeur $8$ et commenter.
\end{enumerate}
\end{exo}

\begin{exo}[Données regroupées en classes \calcoui \corrige]{2}
Une enquête donne les effectifs $4$, $10$, $18$ et $8$ pour les classes
$[0;5[$, $[5;10[$, $[10;15[$ et $[15;20[$.
\begin{enumerate}
  \item Calculer les fréquences et les effectifs cumulés croissants.
  \item Identifier la classe modale et la classe médiane.
  \item Estimer la moyenne à l'aide des centres de classes.
  \item Estimer la variance et l'écart-type.
  \item Expliquer la limite de ces estimations.
\end{enumerate}
\end{exo}

\begin{exo}[Comparer deux distributions \calcoui \corrige]{3}
Deux classes ont obtenu les résultats suivants :
\[
\begin{array}{c|ccc}
 & \text{moyenne} & \text{médiane} & \text{écart-type}\\ \hline
\text{Classe A} & 11{,}8 & 12 & 2{,}1\\
\text{Classe B} & 11{,}8 & 11 & 4{,}3
\end{array}
\]
\begin{enumerate}
  \item Comparer les niveaux moyens des deux classes.
  \item Comparer la dispersion des résultats.
  \item Expliquer ce que suggère la différence entre les médianes.
  \item Rédiger une conclusion de trois phrases sans dépasser les informations
  fournies par le tableau.
\end{enumerate}
\end{exo}

\section{Problèmes de synthèse}

\begin{probleme}[Réservoir et algorithme \calcnon \corrigeP]{2}
Un réservoir contient initialement $120$ litres. Une pompe y ajoute $6$ litres
par minute. Sa capacité maximale est de $360$ litres.
\begin{enumerate}
  \item Calculer les volumes présents après une minute puis après deux minutes.
  \item Exprimer le volume $V(n)$ après $n$ minutes.
  \item Écrire un algorithme qui saisit $n$ et affiche $V(n)$.
  \item Déterminer le premier entier $n$ tel que $V(n)>300$.
  \item Déterminer à partir de quelle minute le réservoir déborde.
\end{enumerate}
\end{probleme}

\begin{probleme}[Tarifs de transport \calcnon \corrigeP]{2}
Une entreprise propose deux tarifs pour un trajet de $x$ kilomètres :
\[
 A(x)=2500+600x
 \qquad\text{et}\qquad
 B(x)=4000+450x,
\]
les montants étant exprimés en ariary.
\begin{enumerate}
  \item Calculer les deux prix pour $x=5$ puis pour $x=20$.
  \item Résoudre l'équation $A(x)=B(x)$.
  \item Déterminer, selon la distance, le tarif le moins cher.
  \item Représenter $A$ et $B$ dans un même repère et interpréter le point
  d'intersection.
  \item Écrire un algorithme qui affiche le tarif le plus avantageux pour une
  distance saisie.
\end{enumerate}
\end{probleme}

\begin{probleme}[Cercle, droite et interprétation \calcnon \corrigeP]{2}
Dans un repère orthonormé, on donne $A(-1;2)$, $B(3;4)$,
$d:x+2y-5=0$ et
\[
 \mathcal C:(x-1)^2+(y-1)^2=9.
\]
\begin{enumerate}
  \item Calculer $\vv{AB}$ et déterminer une équation de $(AB)$.
  \item Déterminer les coordonnées du point d'intersection de $(AB)$ et $d$.
  \item Donner le centre et le rayon de $\mathcal C$.
  \item Étudier la position de $B$ par rapport à $\mathcal C$.
  \item Déterminer les points d'intersection de $d$ et $\mathcal C$, ou
  justifier qu'il n'en existe pas.
\end{enumerate}
\end{probleme}

\begin{probleme}[Terrain rectangulaire \calcnon \corrigeP]{2}
Un terrain rectangulaire a un périmètre de $80$ m. On note $x$ sa largeur.
\begin{enumerate}
  \item Exprimer sa longueur en fonction de $x$.
  \item Justifier que $0<x<40$.
  \item Montrer que son aire est $A(x)=-x^2+40x$.
  \item Mettre $A$ sous forme canonique.
  \item Déterminer les dimensions donnant l'aire maximale et calculer cette aire.
  \item Déterminer les dimensions possibles lorsque l'aire vaut $300$ m$^2$.
\end{enumerate}
\end{probleme}

\begin{probleme}[Boîte sans couvercle \calcoui \corrigeP]{3}
On découpe dans les quatre coins d'une feuille rectangulaire de
$40$ cm sur $20$ cm des carrés de côté entier $x$ cm, puis on relève les bords.
\begin{enumerate}
  \item Justifier que $x$ est un entier vérifiant $1\leq x\leq9$.
  \item Montrer que le volume de la boîte est
  \[
    V(x)=x(40-2x)(20-2x).
  \]
  \item Construire un tableau donnant $V(x)$ pour toutes les valeurs possibles.
  \item Déterminer la valeur de $x$ donnant le plus grand volume dans ce modèle.
  \item Écrire un algorithme qui effectue cette recherche.
  \item Expliquer pourquoi cette étude ne donne pas le maximum parmi toutes
  les valeurs réelles possibles de $x$.
\end{enumerate}
\end{probleme}

\begin{probleme}[Organisation d'une enquête \calcoui \corrigeP]{2}
On a relevé le temps de trajet de $40$ élèves. Les effectifs sont $4$, $10$,
$18$ et $8$ dans les classes $[0;5[$, $[5;10[$, $[10;15[$ et $[15;20[$,
les durées étant exprimées en minutes.
\begin{enumerate}
  \item Construire un tableau contenant les classes, les centres, les effectifs,
  les fréquences et les effectifs cumulés.
  \item Tracer un histogramme de la distribution.
  \item Déterminer la classe modale et la classe médiane.
  \item Estimer la moyenne, la variance et l'écart-type.
  \item Rédiger une conclusion présentant une tendance centrale, une mesure de
  dispersion et la limite liée au regroupement en classes.
\end{enumerate}
\end{probleme}

\begin{probleme}[Trajectoire et sécurité \calcoui \corrigeP]{3}
Dans un repère, la hauteur d'un objet est modélisée par
\[
  h(x)=-x^2+6x+7,
\]
où $x\geq0$ représente la distance horizontale et $h(x)$ la hauteur en mètres.
\begin{enumerate}
  \item Mettre $h$ sous forme canonique.
  \item Déterminer la hauteur maximale et la distance à laquelle elle est atteinte.
  \item Résoudre $h(x)=0$ et interpréter la solution positive.
  \item Un obstacle de hauteur $12$ m se trouve à $x=2$ m. L'objet le franchit-il ?
  \item Déterminer les distances pour lesquelles la hauteur est au moins $12$ m.
\end{enumerate}
\end{probleme}

\begin{probleme}[Cercle circonscrit \calcnon \corrigeP]{3}
Dans un repère orthonormé, on donne $A(-2;0)$, $B(4;2)$ et $C(1;5)$.
\begin{enumerate}
  \item Déterminer le milieu $I$ de $[AB]$ et une équation de la médiatrice
  de $[AB]$.
  \item Déterminer le milieu $J$ de $[AC]$ et une équation de la médiatrice
  de $[AC]$.
  \item Calculer les coordonnées du point $\Omega$ d'intersection des deux
  médiatrices.
  \item Vérifier que $\Omega A=\Omega B=\Omega C$.
  \item Écrire une équation du cercle circonscrit au triangle $ABC$.
\end{enumerate}
\end{probleme}

\begin{tcolorbox}[enhanced,breakable,frame hidden,
  borderline west={2pt}{0pt}{VE},colback=VE!3,sharp corners,
  left=12pt,right=10pt,top=7pt,bottom=7pt]
\textbf{Composition d'une évaluation.} Sélectionner des exercices portant sur
les capacités effectivement enseignées. Une évaluation formative peut réunir
des exercices ciblés et indépendants. Une évaluation sommative peut associer
plusieurs domaines et un problème de synthèse. Adopter une même règle de
calculatrice pour l'ensemble du sujet.
\end{tcolorbox}
